\documentclass{article}
\usepackage{graphicx}
\usepackage{amsmath}
\usepackage{amssymb}
\usepackage{booktabs}
\usepackage{hyperref}
\usepackage[margin=1in]{geometry}
\usepackage{subcaption}
\usepackage{xcolor}

\title{Title}
\date{February 2026}

\begin{document}

\maketitle

\section{Introduction}

%==============================================================================
\section{Experiment 1: Chunk Size Analysis}
%==============================================================================

The purpose of this experiment to identify a minimum text length after which the KL 

\subsection{Motivation}

In order to to generate meaningful texts from llm we need to know how much text do we need for reliable analysis? If it is too small  we have unstable statistics. In this experiment we want study the relationship between chunk size and  stability.

\subsection{Method}

For each author's full book, we:

\begin{enumerate}
    \item \textbf{Extract punctuation sequences}: Parse the text to obtain a sequence of punctuation marks: PUNCTUATION\_VECTOR = ['!', '"', '(', ')', ',', '.', ':', ';', '?', '\^{}']

    
    \item \textbf{Sample chunk pairs}: For each chunk size $s \in \{500, 1000, 1500, 2000, 3000\}$ punctuation marks:
    \begin{itemize}
        \item Sample $n = 30$ pairs of non-overlapping chunks from within the same book
        \item Compute KL divergence between each pair for both f1 and f3
        \item Record mean and standard deviation: $\mu \pm \sigma$
    \end{itemize}
    
    \item \textbf{Visual comparison}: For each chunk size, extract:
    \begin{itemize}
        \item Chunk A: From the midpoint of the \emph{first half} of the book
        \item Chunk B: From the midpoint of the \emph{second half} of the book
    \end{itemize}
    These positions ensure the chunks are well-separated (avoiding overlap) while representing typical content (avoiding beginning/end effects).
\end{enumerate}

\subsection{Results}

\subsubsection{Heatmap Visualisation}

Figure~\ref{fig:heatmaps} shows the punctuation sequences as coloured strips, where each colour represents a different punctuation mark. Each row shows one author; each column pair shows chunks from the first and second half of the book at a given chunk size.

\textbf{Key observations}:
\begin{itemize}
    \item At 500 marks, patterns are difficult to discern visually
    \item At 2000+ marks, stylistic patterns become apparent (e.g., comma-period dominance)
\end{itemize}

\begin{figure}[htbp]
    \centering
    \includegraphics[width=\textwidth]{plot1_heatmaps_comparison.png}
    \caption{Punctuation sequence heatmaps at different chunk sizes. Each pair shows chunks from the midpoint of the first half vs.\ the midpoint of the second half of the same book. Visual consistency within author (across columns) indicates  stability.}
    \label{fig:heatmaps}
\end{figure}

\subsubsection{KL Divergence Stability}

Figure~\ref{fig:kl_variance} shows how the within-book KL divergence varies with chunk size for each author. The shaded regions represent $\pm 1$ standard deviation.

\textbf{Key observations}:
\begin{itemize}
    \item \textbf{Variance decreases} as chunk size increases --- larger samples give more stable estimates
    \item \textbf{Mean KL stabilises} around 1500--2000 marks for most authors
\end{itemize}

\begin{figure}[htbp]
    \centering
    \includegraphics[width=\textwidth]{plot2_kl_variance.png}
    \caption{KL divergence (mean $\pm$ std.\ dev.) for within-book chunk comparisons. Blue: f1 (frequency); Red: f3 (joint probability). Lower values indicate more consistent punctuation patterns.}
    \label{fig:kl_variance}
\end{figure}

%==============================================================================
\section{Experiment 2: Baseline Comparison}
%==============================================================================

What is a sensible baseline? 

Sam Howison
The key thing will be to pin down exactly what will be the best questions. If we ask 'how well do LLMs match the actual punctuation sequences' it would be good to have a null model as well. Eg an average over contemporary authors. ****The we can say things like 'LLM writing Austen is much closer (or not)  to Austen than a random contemporary (whatever) author (+ quantify via K-L).**** Does that make sense as a starting question??


PA
Take chunks of 2,000 punctuation marks from 50 different authors
Concatenate them into one synthetic "book"
Compute within-book KL on this synthetic book = baseline
Compare each real author's within-book KL to this baseline
\begin{itemize}
\item first define "random author": two options (one is what you have, the other is literally <f1> over 50 authors)
\item second measure different things: "is LLM writing austen more similar to austen than LLM writing austen is to random author" or/and "is LLM writing austen more similar to austen than austen is to random author" 
\end{itemize}


MB
KL(author-f1, LLM-f1) 

with 

KL(author-f1, baseline-f1) 

Is the first "more similar"? 

potential candidate for baseline f1: 

sum the f1 of N distinct authors chosen at random then divide by N (say N = 50). This gives an "average/baseline f1" 

Not sure if that makes sense?

I guess we can also control for length of docs from which f1 is computed if we want to so LLM sequence = author sequence = minimum we identified as meaningful

%==============================================================================
\section{Experiment 3: LLM comparison}
%==============================================================================

\subsection{Setup}

For each of the four author i generated a 2000 mark piece of text (10X for robustness checks)

\begin{table}[htbp]
\centering
\begin{tabular}{lccc}
\toprule
Author & Real marks & Generated marks (range) & Chunk size \\
\midrule
Jane Austen & 25,935 & 2,108 -- 2,853 & 2,000 \\
William Shakespeare & 6,253 & 2,006 -- 2,807 & 2,000 \\
H.G. Wells & 10,099 & 2,073 -- 2,432 & 2,000 \\
Agnes May Fleming & 18,922 & 2,000 -- 2,578 & 2,000 \\
\bottomrule
\end{tabular}
\caption{Overview of real and generated text lengths for each author.}
\label{tab:setup}
\end{table}

\subsection{Baseline 1: LLM-to-target vs LLM-to-random}

For each of the 10 LLM runs
\begin{itemize}
    \item we compute $d_{KL}$(LLM, Target) and $d_{KL}$(LLM, Random$_i$) for all 47 random Gutenberg authors.
    \item If the LLM is actually imitating the author, the LLM-to-author distance should be lower than LLM-to-random.
\end{itemize}

\subsubsection{f1 (frequency distribution)}

\begin{table}[htbp]
\centering
\begin{tabular}{lcc}
\toprule
Author & KL(LLM, Target) & KL(LLM, Random) \\
\midrule
Jane Austen & $0.120 \pm 0.041$ & $0.178 \pm 0.172$ \\
Shakespeare & $0.094 \pm 0.052$ & $0.264 \pm 0.127$ \\
H.G. Wells & $0.139 \pm 0.054$ & $0.200 \pm 0.204$ \\
A.M. Fleming & $0.147 \pm 0.034$ & $0.177 \pm 0.155$ \\
\bottomrule
\end{tabular}
\caption{Baseline 1 results for f1 (frequency distribution).}
\label{tab:baseline1_f1}
\end{table}

\subsubsection{f3 (transition matrix)}

\begin{table}[htbp]
\centering
\begin{tabular}{lcc}
\toprule
Author & KL(LLM, Target) & KL(LLM, Random) \\
\midrule
Jane Austen & $0.263 \pm 0.072$ & $0.349 \pm 0.261$ \\
Shakespeare & $0.201 \pm 0.107$ & $0.512 \pm 0.265$ \\
H.G. Wells & $0.283 \pm 0.100$ & $0.354 \pm 0.284$ \\
A.M. Fleming & $0.279 \pm 0.055$ & $0.330 \pm 0.230$ \\
\bottomrule
\end{tabular}
\caption{Baseline 1 results for f3 (transition matrix).}
\label{tab:baseline1_f3}
\end{table}

\subsection{Baseline 2: LLM's distance to Original Author vs random authors' distance to Original Author}

Compare $d_{KL}$(LLM\_run, Real\_Target) vs $d_{KL}$(Random$_i$, Real\_Target). If LLM-to-author < Random-to-target, the LLM better at matching the target author than a random author

\subsubsection{f1 (frequency distribution)}

\begin{table}[htbp]
\centering
\begin{tabular}{lcc}
\toprule
Author & KL(LLM, Target) & KL(Random, Target) \\
\midrule
Jane Austen & $0.120 \pm 0.041$ & $0.158 \pm 0.149$ \\
Shakespeare & $0.094 \pm 0.052$ & $0.112 \pm 0.092$ \\
H.G. Wells & $0.139 \pm 0.054$ & $0.314 \pm 0.231$ \\
A.M. Fleming & $0.147 \pm 0.034$ & $0.176 \pm 0.163$ \\
\bottomrule
\end{tabular}
\caption{Baseline 2 results for f1 (frequency distribution).}
\label{tab:baseline2_f1}
\end{table}

\subsubsection{f3 (transition matrix)}

\begin{table}[htbp]
\centering
\begin{tabular}{lcc}
\toprule
Author & KL(LLM, Target) & KL(Random, Target) \\
\midrule
Jane Austen & $0.263 \pm 0.072$ & $0.279 \pm 0.187$ \\
Shakespeare & $0.201 \pm 0.107$ & $0.274 \pm 0.174$ \\
H.G. Wells & $0.283 \pm 0.100$ & $0.637 \pm 0.395$ \\
A.M. Fleming & $0.279 \pm 0.055$ & $0.389 \pm 0.313$ \\
\bottomrule
\end{tabular}
\caption{Baseline 2 results for f3 (transition matrix).}
\label{tab:baseline2_f3}
\end{table}

\subsection{heatmaps visualisation}

However, the heatmaps and frequency plots do not visually confirm this\ldots

Figure~\ref{fig:llm_heatmaps} shows punctuation sequences as coloured strips for both real and LLM-generated texts. Each row represents a different author, with the real text shown on the left and LLM-generated text on the right.

\textbf{Key observations}:
\begin{itemize}
    \item Real texts show rich, varied colour patterns
    \item LLM texts are dominated by comma (light blue) and period (yellow)
    \item Low diversity in LLM punctuation patterns across all authors
\end{itemize}

\begin{figure}[htbp]
    \centering
    \includegraphics[width=\textwidth]{all_authors_heatmaps.png}
    \caption{Punctuation sequence heatmaps comparing real author texts (left) vs LLM-generated texts (right) for 2000-mark chunks. The visual disparity reveals that despite passing relative statistical tests, LLM texts exhibit a generic, comma-heavy punctuation fingerprint.}
    \label{fig:llm_heatmaps}
\end{figure}

The baselines above are relative comparisons (is the LLM closer to the target than to random?). They pass, meaning there is a detectable signal of style adaptation. However, the heatmaps and frequency plots reveal that in absolute terms, the imitation is shallow. Taking Austen as a representative example in figure~\ref{tab:austen_comparison}

\begin{table}[htbp]
\centering
\begin{tabular}{lcc}
\toprule
Punctuation mark & Real Austen & LLM-Austen \\
\midrule
Comma & 38.3\% & $\sim$58\% \\
Period & 27.6\% & $\sim$27\% \\
Quote marks & 19.3\% & $\sim$10\% \\
Semicolon & 8.6\% & $\sim$0.3\% \\
Exclamation & 3.3\% & $\sim$0.5\% \\
Colon & 0.5\% & $\sim$0\% \\
\bottomrule
\end{tabular}
\caption{Comparison of punctuation frequencies: Real Austen vs LLM-generated Austen.}
\label{tab:austen_comparison}
\end{table}

Real texts show rich, varied colour patterns, while LLM texts are dominated by comma (light blue) and period (yellow). The LLM has a generic punctuation fingerprint comma-and-period heavy, with low diversity.

%==============================================================================
\section{Experiment 4: Prompt Strategy Comparison}
%==============================================================================

\subsection{Motivation}

The results of Experiment~3 show that the LLM (Gemini 2.5 Flash) produces text that is closer to the target author than to random authors in relative terms, but the heatmaps reveal the imitation is shallow in absolute terms. A natural question is: \textbf{can we improve the LLM's punctuation fidelity by changing the prompt?}

We test three prompting strategies, each representing a different level of information given to the model:

\begin{description}
    \item[v1 -- Excerpt only (baseline):] The prompt includes a 1500-word excerpt from the middle of the author's real book, followed by the instruction ``Write an original piece of fiction in the same writing style as the passage above.'' No explicit mention of punctuation.
    \item[v2 -- Excerpt + explicit punctuation targets:] The same excerpt as v1, but augmented with a detailed punctuation style guide computed from the real book. This guide includes: (i)~exact punctuation frequency percentages (e.g.\ ``comma: 45.2\%, period: 28.1\%''), (ii)~transition probabilities (e.g.\ ``after a comma, the next punctuation is typically: comma 42\%, period 28\%''), (iii)~sentence length distribution (mean, quartiles, bucket percentages), and (iv)~punctuation density (average words between marks). The LLM is instructed to follow these targets closely.
    \item[v3 -- Author name only:] No excerpt and no statistics. The prompt simply says ``Write an original piece of fiction in the writing style of [Author Name].'' This tests the LLM's parametric (training-time) knowledge of each author's style.
\end{description}

All other settings are identical: Gemini 2.5 Flash, temperature 1.0, target of 2000 punctuation marks, em-dash post-processing. Each version generates 10 runs per author. All results are measured using the same baseline methodology from Experiment~3 (2000-mark chunks, same 47 random Gutenberg authors).

\subsection{Reproducibility check}

Before comparing versions, we checked whether v1 results are reproducible. We re-ran v1 for Jane Austen twice (10 runs each) and applied the baseline analysis to all three batches.

\begin{table}[htbp]
\centering
\begin{tabular}{lcccc}
\toprule
Batch & $d_{KL}$(LLM, Target) f1 & $d_{KL}$(LLM, Target) f3 \\
\midrule
v1 original  & $0.120 \pm 0.041$ & $0.263 \pm 0.072$ \\
v1 rerun 1   & $0.116 \pm 0.042$ & $0.247 \pm 0.069$ \\
v1 rerun 2   & $0.119 \pm 0.028$ & $0.254 \pm 0.053$ \\
\bottomrule
\end{tabular}
\caption{Reproducibility of v1 baseline for Jane Austen (3 independent batches of 10 runs each, all measured via the standardised baseline analysis). The means are stable across batches ($\sim$0.12 for f1, $\sim$0.25 for f3).}
\label{tab:reproducibility}
\end{table}

Note: the generation script's inline KL computation (which reported higher values of 0.23 and 0.27 for the reruns) uses a different chunking strategy than the baseline analysis. The baseline analysis, which uses consistent 2000-mark middle chunks from the real text, gives stable and reproducible results.

\subsection{Results: Q2 comparison across prompt strategies}

We report Q2 results ($d_{KL}$(LLM, Target) vs $d_{KL}$(Random, Target)) for all four authors across all three prompt versions. Bold indicates the LLM is closer to the target than random (i.e.\ the LLM ``beats'' the baseline).

\subsubsection{f1 (frequency distribution)}

\begin{table}[htbp]
\centering
\begin{tabular}{lcccc}
\toprule
Author & v1 (excerpt) & v2 (excerpt+targets) & v3 (name only) & Random \\
\midrule
Austen   & $\mathbf{0.120}$ & $\mathbf{0.115}$ & $\mathbf{0.149}$ & $0.158$ \\
Shakespeare & $\mathbf{0.094}$ & $0.396$ & $0.122$ & $0.112$ \\
Wells    & $\mathbf{0.139}$ & $\mathbf{0.287}$ & $0.409$ & $0.314$ \\
Fleming  & $\mathbf{0.147}$ & $\mathbf{0.161}$ & $0.203$ & $0.176$ \\
\bottomrule
\end{tabular}
\caption{$d_{KL}$(LLM, Target) for f1 across prompt strategies. Bold = LLM closer to target than random baseline. v1 is the best or tied-best for every author.}
\label{tab:v_comparison_f1}
\end{table}

\subsubsection{f3 (transition matrix)}

\begin{table}[htbp]
\centering
\begin{tabular}{lcccc}
\toprule
Author & v1 (excerpt) & v2 (excerpt+targets) & v3 (name only) & Random \\
\midrule
Austen   & $\mathbf{0.263}$ & $\mathbf{0.272}$ & $0.314$ & $0.279$ \\
Shakespeare & $\mathbf{0.201}$ & $0.768$ & $0.297$ & $0.274$ \\
Wells    & $\mathbf{0.283}$ & $0.654$ & $0.826$ & $0.637$ \\
Fleming  & $\mathbf{0.279}$ & $\mathbf{0.289}$ & $\mathbf{0.340}$ & $0.389$ \\
\bottomrule
\end{tabular}
\caption{$d_{KL}$(LLM, Target) for f3 across prompt strategies. Bold = LLM closer to target than random baseline. v1 is the best for every author on f3.}
\label{tab:v_comparison_f3}
\end{table}

\subsection{Discussion}

\subsubsection{v1 (excerpt only) is the best strategy overall}

For every author and both features, v1 achieves the lowest or tied-lowest KL divergence to the target. The LLM is better at implicitly absorbing punctuation patterns from a passage than at following explicit statistical instructions (v2) or relying on parametric knowledge alone (v3).

\subsubsection{Explicit punctuation targets (v2) hurt rather than help}

Adding numerical punctuation guidelines on top of the excerpt either leaves performance unchanged (Austen, Fleming) or makes it significantly worse (Shakespeare, Wells). The effect is catastrophic for Shakespeare: f1 jumps from 0.094 to 0.396 (3.5$\times$ worse) and f3 from 0.201 to 0.768 (3.8$\times$ worse). Only 8 of 10 Shakespeare runs produced enough punctuation marks, suggesting the numerical constraints also disrupted generation fluency.

This suggests the LLM cannot reliably follow quantitative statistical constraints on low-level textual properties. The explicit targets may conflict with the model's implicit style model, producing stilted text with worse punctuation patterns than unconstrained generation.

\subsubsection{Author name alone (v3) is insufficient}

Without an excerpt, performance degrades for all authors except Shakespeare on f1. For Wells, v3 is the worst condition: Q1 is \emph{negative} (the LLM is closer to random authors than to Wells), meaning the model has no useful internal representation of Wells's punctuation style.

Shakespeare is the exception: v3 achieves $0.122$ on f1, competitive with v1's $0.094$, likely because Shakespeare is heavily represented in the training data. However, even for Shakespeare, v3 fails to beat the random baseline on Q2 for both features.

\subsubsection{The excerpt does the heavy lifting}

The comparison between v1 and v3 isolates the contribution of the in-context excerpt. For Austen (f1: 0.120 vs 0.149), the excerpt provides a modest improvement. For Wells (f1: 0.139 vs 0.409), the excerpt is essential --- without it, the LLM produces punctuation patterns further from the target than random authors.

\begin{table}[htbp]
\centering
\begin{tabular}{lcccc}
\toprule
 & \multicolumn{2}{c}{f1 (frequency)} & \multicolumn{2}{c}{f3 (transition)} \\
\cmidrule(lr){2-3}\cmidrule(lr){4-5}
Author & Beats random? & Best version & Beats random? & Best version \\
\midrule
Austen    & v1, v2, v3 & v2 & v1, v2 & v1 \\
Shakespeare & v1 only & v1 & v1 only & v1 \\
Wells     & v1, v2 & v1 & v1 only & v1 \\
Fleming   & v1, v2 & v1 & v1, v2, v3 & v1 \\
\bottomrule
\end{tabular}
\caption{Summary: which prompt versions beat the random baseline (Q2), and which version achieves the lowest $d_{KL}$ to target.}
\label{tab:v_summary}
\end{table}

%==============================================================================
\section*{References}
%==============================================================================


\end{document}